\chapter{Resultados}
\label{chap:resultados}

En este capítulo se van a mostrar los resultados del juego Retro Fantasy, el alcance de la plataforma y cuál ha sido el nivel de cumplimientos de los objetivos que se establecieron al inicio del desarrollo.

\section{Resultados de la Integración de Datos y Motor de Juego}
\label{sec:resultados_datos}

Uno de los principales logros técnicos del proyecto es el procesamiento e integración del dataset de Kaggle. Los resultados más destacables en este ámbito son:

\begin{itemize}
    \item \textbf{Carga y Migración Exitosa del Catálogo}: Se ha migrado y estructurado un volumen de datos correspondiente a más de 10.000 jugadores y 25.000 partidos históricos de la base de datos SQLite de Kaggle al motor relacional de PostgreSQL en Supabase.
    \item \textbf{Parseo de Micro-Eventos XML en Tiempo Real}: El motor del backend ha demostrado una precisión del 100\% en la lectura y agregación de las etiquetas XML de la tabla \texttt{Match}. Se han procesado con éxito eventos complejos como goles, tarjetas, tiros, etc. Para traducirlos dinámicamente a notas decimales y picas subjetivas por medio del cronista virtual.
    \item \textbf{Consistencia Transaccional}: Las pruebas realizadas confirman la robustez del sistema frente a condiciones de carrera, como pujas duplicadas o clausulazos simultáneos, gracias al uso de bloqueos explícitos (\texttt{FOR UPDATE}) y transacciones SQL atómicas, manteniendo los presupuestos y plantillas siempre consistentes.
    \item \textbf{Conversión y Pujas Simuladas de la Máquina}: Se ha verificado la correcta ejecución del motor de cierre de mercado en presencia de jugadores en venta listados por mánagers, comprobando que se generan pujas automáticas del bot en el rango estipulado ([-10\%, +10\%]) y que todas las pujas (tanto humanas como de la máquina) se convierten correctamente en ofertas directas pendientes. Esto permite que la resolución de la venta sea exclusiva y manual por parte del mánager propietario, liquidando saldos y actualizando plantillas de forma atómica al aceptarse.
\end{itemize}

\section{Resultados de la Interfaz y Vistas de la Aplicación}
\label{sec:resultados_interfaz}

La aplicación web se ha desarrollado implementando una estética de tipo \textit{Liquid Glass} y maquetación con \textit{Bento Cards}, optimizada para dispositivos móviles mediante un diseño pensado en estos. 

Para demostrar la implementación y funcionamiento visual de cada pantalla, a continuación se detallan las principales vistas funcionales de la aplicación, incluyendo los contenedores para las capturas de pantalla de la interfaz final.

\subsection{Pantalla de Inicio de Sesión y Registro}
Esta pantalla de acceso a la plataforma está integrada con Supabase Auth. Garantiza una entrada rápida y segura para el usuario, validando las credenciales y almacenando el token JWT.

\begin{figure}[H]
    \centering
    \includegraphics[width=\textwidth]{Plantilla de Trabajo Fin de Grado - New Version 1/figures/Login.png}
    \caption{Interfaz de acceso y registro de usuarios en Retro Fantasy.}
    \label{fig:res_login}
\end{figure}

\subsection{Dashboard de Mánager (Home)}
Es el centro de control del usuario, que muestra un resumen rápido de las estadísticas más necesarias del juego. Muestra Bento Cards con la jornada activa en juego, los puntos de tu rival más cercano, las mayores subidas y bajadas del mercado.

\begin{figure}[H]
    \centering
    \includegraphics[width=\textwidth]{Plantilla de Trabajo Fin de Grado - New Version 1/figures/Dashboard.png}
    \caption{Dashboard principal del mánager en Retro Fantasy.}
    \label{fig:res_dashboard}
\end{figure}

\subsection{Clasificación General y Equipos Rivales}
Muestra la vista de la clasificación de los participantes de la liga ordenados por su puntuación acumulada. Permite hacer clic en cualquier rival para auditar su plantilla actual y realizar ofertas directas de transferencia o realizar un clausulazo.

\begin{figure}[H]
    \centering
    \includegraphics[width=\textwidth]{Plantilla de Trabajo Fin de Grado - New Version 1/figures/Clasificacion.png}
    \caption{Tabla de clasificación de la liga y acceso a plantillas rivales.}
    \label{fig:res_clasificacion}
\end{figure}

\subsection{Mercado de Fichajes y Pujas}
Muestra los agentes libres transferibles que publica el sistema diariamente. Incluye el valor de mercado base, un temporizador de cuenta atrás para el cierre del ciclo y el formulario para registrar y modificar las ofertas de compra.

\begin{figure}[H]
    \centering
    \includegraphics[width=\textwidth]{Plantilla de Trabajo Fin de Grado - New Version 1/figures/Mercado.png}
    \caption{Mercado de fichajes activo con listado de agentes libres y campos de puja.}
    \label{fig:res_mercado}
\end{figure}

\subsection{Vista de Alineación y Gestión de Plantilla}
Representa de forma interactiva el terreno de juego. Permite arrastrar los jugadores de la plantilla hacia el once titular, modificar la formación (ej. 4-4-2, 3-5-2) y gestionar los blindajes de las cláusulas de rescisión de cada futbolista.

\begin{figure}[H]
    \centering
    \includegraphics[width=\textwidth]{Plantilla de Trabajo Fin de Grado - New Version 1/figures/Alineacion.png}
    \caption{Vista interactiva del campo de fútbol para la alineación y tácticas.}
    \label{fig:res_alineacion}
\end{figure}

\section{Referencias a los Anexos y Guías Visuales}
\label{sec:res_anexos}

Para no sobrecargar el cuerpo principal de la memoria con un número excesivo de capturas de flujos de navegación secundarios (como los formularios de perfil, la creación de tickets de soporte o la recuperación de contraseñas), el manual detallado de la aplicación y las guías paso a paso de uso se han recopilado e integrado en el \textbf{Apéndice~\ref{app:manual_usuario}} (Manual de Usuario). De este modo, se garantiza una lectura más fluida de la memoria de diseño sin descuidar el detalle documental.
